\chapterhead{57}{ORR-18}{Risk Capital Attribution and Risk-Adjusted Performance Measurement}{1}{9}

\lohead{57}{a}{Define, compare, and contrast risk capital, economic capital, and
regulatory capital, and explain methods and motivations for using economic capital
approaches to allocate risk capital.}

\orsub{Risk capital, economic capital, and regulatory capital}

\noindent\begin{minipage}{\textwidth}
\renewcommand{\arraystretch}{1.35}
\begin{tabularx}{\textwidth}{@{}L{0.24\textwidth}Y@{}}
\rowcolor{navy} \thd{Type} & \thd{Definition}\\
{\tabfont \textbf{1) Risk capital}} &
{\tabfont Provides a financial buffer to protect against unexpected losses.}\\
\rowcolor{paleblue}
{\tabfont \textbf{2) Economic capital}} &
{\tabfont The amount of capital a firm needs to sustain risks, both expected and
unexpected, and continue as a going concern.}\\
{\tabfont \textbf{3) Regulatory capital}} &
{\tabfont Mandated by regulatory bodies specifically for regulated industries like
banking and insurance. Regulatory capital is specific to regulated sectors.}\\
\end{tabularx}
\end{minipage}
\orfigcap{Three capital concepts, distinguished.}

\orsub{Using economic capital approaches}

\begin{lettered}
\item It cushions the risks arising from leverage and financial activities like
derivatives.
\item It assures creditors (depositors) of the institution's creditworthiness.
\item It acts as a buffer against unexpected events and helps maintain financial
stability. However, holding economic capital increases the cost of capital and
reduces profits.
\item Finding the right balance between holding enough capital and taking risks is
crucial.
\end{lettered}

% ---------------------------------------------------------------------
\lohead{57}{b}{Describe the RAROC (risk-adjusted return on capital) methodology and
its use in capital budgeting.}

\orsub{Risk-adjusted return on capital}

RAROC measures risk-adjusted performance, not just raw numbers. It allocates risk
capital for a more accurate view.

\orsubb{Benefits}

\begin{lettered}
\item Performance measurement using economic profits instead of accounting
profits.
\item Links to shareholder value creation for bonuses and performance evaluations.
\item Helps decide on buying or selling businesses, or starting or stopping
projects, based on risk-adjusted value.
\item Using risk-based pricing, which allows proper pricing that takes into account
the economic risks undertaken by a firm in a given transaction.
\end{lettered}

% ---------------------------------------------------------------------
\lohead{57}{c}{Compute and interpret the RAROC for a project, loan, or loan
portfolio and use RAROC to compare business unit performance.}

\orsub{Two formulas to calculate RAROC}

\begin{orkeybox}
\[
\mathrm{RAROC}=\frac{\text{after-tax expected risk-adjusted net income}}
                    {\text{economic capital}}
\]
\vspace{2pt}
\[
\mathrm{RAROC}=\frac{\left(\begin{array}{c}
\text{expected revenues}-\text{costs}-\text{expected losses}\\
-\ \text{taxes}+\text{return on economic capital}\pm\text{transfers}
\end{array}\right)}{\text{economic capital}}
\]
\end{orkeybox}

\begin{ornotebox}{Note}
RAROC is similar to the Sharpe ratio and NPV, but considers both systematic and
unsystematic risk. It uses economic capital instead of accounting capital, and
informs decisions on projects, pricing, and capital allocation.
\end{ornotebox}

\begin{ornotebox}{Remember}
\begin{enumerate}
\item[\textcolor{gold}{\sffamily\bfseries 1)}] For \textbf{capital budgeting
projects}, use \textbf{expected} revenues and losses.
\item[\textcolor{gold}{\sffamily\bfseries 2)}] For \textbf{performance
evaluation}, use \textbf{actual} revenues and losses.
\end{enumerate}
\end{ornotebox}

\begin{orexambox}{Example: RAROC calculation}
Assume the following information for a commercial loan portfolio:
\begin{itemize}
\item \$1.5 billion principal amount
\item 7\% pre-tax expected return on loan portfolio
\item Direct annual operating costs of \$10 million
\item Loan portfolio is funded by \$1.5 billion of retail deposits; interest rate
is 5\%
\item Expected loss on the portfolio is 0.5\% of principal per annum
\item Unexpected loss of 8\% of the principal amount, or \$120 million of economic
capital required
\item Risk-free rate on government securities is 1\% (based on the economic
capital required)
\item 25\% effective tax rate
\item Assume no transfer pricing issues
\end{itemize}
\textbf{Compute} the RAROC for this loan portfolio.

\vspace{7pt}
\noindent{\sffamily\bfseries\fontsize{8.6}{10}\selectfont\color{forest}Answer}\par
\vspace{3pt}
First, calculate the following RAROC components:
\begin{itemize}
\item Expected revenue $= 0.07 \times \$1.5\text{ billion} = \$105$ million
\item Interest expense $= 0.05 \times \$1.5\text{ billion} = \$75$ million
\item Expected loss $= 0.005 \times \$1.5\text{ billion} = \$7.5$ million
\item Return on economic capital $= 0.01 \times \$120\text{ million} = \$1.2$
million
\end{itemize}
Then, apply the RAROC equation:
\[
\mathrm{RAROC}=\frac{(105-10-75-7.5+1.2+0)\times(1-0.25)}{120}=8.56\%
\]
Therefore, maintenance of the commercial loan portfolio requires an after-tax
expected rate of return on equity of at least 8.56\%.
\end{orexambox}

% ---------------------------------------------------------------------
\lohead{57}{d}{Explain challenges that arise when using RAROC for performance
measurement, including choosing a time horizon, measuring default probability, and
choosing a confidence level.}

\orsubb{1) Time horizon}

Typically, RAROC is calculated using a one-year horizon to align with business
planning cycles and risk recovery capabilities. However, longer-term transactions
may necessitate a multi-period RAROC to capture changes in risk over time, which
can lead to periods of over- or under-capitalisation if using an averaging
approach.

\orsubb{2) Default probability}

\noindent\begin{minipage}{\textwidth}
\renewcommand{\arraystretch}{1.35}
\begin{tabularx}{\textwidth}{@{}L{0.26\textwidth}Y@{}}
\rowcolor{navy} \thd{Approach} & \thd{Use}\\
{\tabfont \textbf{Point-in-Time (PIT)}} &
{\tabfont Used for short-term loss expectations and pricing; sensitive to current
conditions.}\\
\rowcolor{paleblue}
{\tabfont \textbf{Through-the-Cycle (TTC)}} &
{\tabfont More stable, and used for long-term economic capital assessments and
strategic planning. It results in lower volatility of economic capital compared to
the PIT approach.}\\
\end{tabularx}
\end{minipage}
\orfigcap{PIT reacts to conditions; TTC smooths through them.}

\orsubb{3) Confidence level}

\begin{itemize}
\item Higher credit ratings require higher confidence levels for capital
allocation.
\item Lower confidence levels reduce capital needed but impact performance
metrics.
\end{itemize}

% ---------------------------------------------------------------------
\lohead{57}{e}{Calculate the hurdle rate and apply this rate in making business
decisions using RAROC.}

\orsub{Hurdle rate for capital budgeting decisions}

Most firms use a single hurdle rate for all business activities: the after-tax
weighted-average cost of equity capital.

\begin{orkeybox}
\[
h_{AT}=\frac{\bigl[(CE \times R_{CE}) + (PE \times R_{PE})\bigr]}{(CE+PE)}
\]
\end{orkeybox}

\begin{itemize}
\item The cost of preferred equity, $r_{PE}$, is simply the yield on the firm's
preferred shares.
\item The cost of common equity, $r_{CE}$, is determined via a model such as the
CAPM:
\[
R_{CE}=R_{F}+\beta_{CE}\,(R_{M}-R_{F})
\]
\end{itemize}

\begin{orexambox}{Decision rule}
If RAROC $>$ hurdle rate, \textbf{accept} the project. Otherwise, reject.
\end{orexambox}

% ---------------------------------------------------------------------
\lohead{57}{f}{Compute the adjusted RAROC for a project to determine its viability.}

\orsub{Adjusted RAROC}

Adjusting the traditional RAROC calculation obtains a RAROC measure that takes into
account the systemic riskiness of returns, and for which the hurdle rate is the
same across all business lines.

\begin{orkeybox}
\[
\text{Adjusted RAROC} = \mathrm{RAROC} - \beta_{E}\,(R_{M}-R_{F})
\]
\end{orkeybox}

\begin{orexambox}{New decision rule}
If adjusted RAROC $> R_{F}$, \textbf{accept} the project. Otherwise, reject.
\end{orexambox}

\begin{orexambox}{Example: adjusted RAROC}
Suppose RAROC is 12\%, the risk-free rate is 5\%, the market return is 11\%, and
the firm's equity beta is 1.5. Use ARAROC to \textbf{determine} whether the
project should be accepted or rejected.

\vspace{7pt}
\noindent{\sffamily\bfseries\fontsize{8.6}{10}\selectfont\color{forest}Answer}\par
\vspace{3pt}
\[
\text{Adjusted RAROC} = \mathrm{RAROC} - \beta_{E}\,(R_{M}-R_{F})
= 0.12 - 1.5(0.11-0.05) = 0.12-0.09 = 0.03
\]
The project should be rejected because the ARAROC of 3\% is less than the
risk-free rate of 5\%.
\end{orexambox}

% ---------------------------------------------------------------------
\lohead{57}{g}{Explain challenges in modeling diversification benefits, including
aggregating a firm's risk capital and allocating economic capital to different
business lines.}

\orsub{Risk capital and diversification}

\begin{orkeybox}
The overall risk capital for a firm should be \emph{less} than the total of the
individual risk capitals of the underlying business units. This is because the
correlation of returns between business units is likely to be less than $+1$.
\end{orkeybox}

A business unit with earnings or cash flows that are highly correlated to the
overall firm should be allocated more risk capital than a business unit with
earnings or cash flows that are negatively correlated, assuming similar
volatility. Having business lines that are countercyclical in nature allows the
overall firm to have stable earnings and to attain a given desired credit rating
using less risk capital.

% ---------------------------------------------------------------------
\lohead{57}{h}{Explain best practices in implementing an approach that uses RAROC
to allocate economic capital.}

\orsub{RAROC best practices}

Implementing a RAROC approach effectively involves several key strategies.

\begin{lettered}
\item \textbf{Senior management involvement.} The management team, including the
CEO, should actively promote RAROC as a key metric for evaluating shareholder value
creation and profit relative to risk.
\item \textbf{Communication and education.} Ensure all management levels
understand RAROC processes to secure buy-in, especially regarding how economic
capital is allocated across business units.
\item \textbf{Ongoing consultation.} Establish a committee to periodically review
and update key metrics related to credit, market, and operational risks to ensure
fair and accurate capital allocation.
\item \textbf{Data quality control.} Centralise data collection with checks for
accuracy and reasonability to support reliable RAROC calculations.
\item \textbf{Complement RAROC with qualitative factors.} Assess business units not
just on RAROC returns but also on the quality of earnings and strategic importance,
using a four-quadrant analysis to guide resource allocation.
\item \textbf{Active capital management.} Business units should regularly submit
their capital and risk limit requests for evaluation and adjustment by the RAROC
team and senior management, ensuring alignment with overall strategic and funding
limits.
\end{lettered}
